% 本科目录格式：章节号前用\quad\quad缩进（2em），与原始模板一致
\ifcauc@bachelor
\titlecontents{chapter}[0pt]{\heiti\zihao{-4}}{\thecontentslabel\quad}{}
        {\titlerule*[10pt]{$\ldots$}\contentspage}
\titlecontents{section}[0pt]{\songti\zihao{-4}}{\quad\quad\thecontentslabel\quad}{}
        {\titlerule*[10pt]{$\ldots$}\contentspage}
\else
% 研究生目录格式由下方 \dottedcontents 定义
% 原始模板规范：研究生目录section使用\quad\quad前缀（与本科一致），subsection不单独定义
\titlecontents{chapter}[0pt]{\heiti\zihao{-4}}{\thecontentslabel\quad}{}
        {\titlerule*[10pt]{$\ldots$}\contentspage}
\titlecontents{section}[0pt]{\songti\zihao{-4}}{\quad\quad\thecontentslabel\quad}{}
        {\titlerule*[10pt]{$\ldots$}\contentspage}
\fi

% 本科目录使用阿拉伯数字章节号+点引导线，研究生目录使用中文数字章节号，符合各学位论文格式规范
\ifcauc@bachelor
\ctexset{
    chapter = {
        nameformat = {\bfseries\zihao{3}},
        format={\bfseries\zihao{3}},
        number={\arabic{chapter}},
        aftername={\quad},
        name={第~,~章},
        beforeskip=-5pt,
        afterskip=30pt,
        pagestyle=fancy
    },
    section = {
        name={},
        format={\bfseries\zihao{4}},
        aftername={\quad},
        beforeskip=7pt,
        afterskip=7pt
    },
    subsection = {
        name={},
        format={\heiti\zihao{-4}},
        aftername={\quad},
        beforeskip=7pt,
        afterskip=7pt
    },
    % 原始模板规范：本科subsubsection间距7pt/7pt
    subsubsection = {
        name={},
        format={\heiti\zihao{-4}},
        aftername={\quad},
        beforeskip=7pt,
        afterskip=7pt
    },
    contentsname={\centerline{\heiti\zihao{3}\texorpdfstring{目~~录}{目录}}},
    bibname={\centerline{\heiti\zihao{3}参考文献}},
    paragraph={format={\zihao{-4}}}
}
\else
\ctexset{
    chapter = {
        nameformat = {\heiti\zihao{3}},
        format={\heiti\zihao{3}},
        number={\zhnum{chapter}},
        aftername={\enspace},
        name={第,章},
        beforeskip=-5pt,
        afterskip=30pt,
        pagestyle=fancy
    },
    section = {
        name={},
        format={\bfseries\zihao{4}},
        aftername={\enspace},
        beforeskip=24pt,
        afterskip=6pt
    },
    subsection = {
        name={},
        format={\bfseries\zihao{4}},
        aftername={\enspace},
        beforeskip=12pt,
        afterskip=6pt
    },
    subsubsection = {
        name={},
        format={\bfseries\zihao{4}},
        aftername={\enspace},
        beforeskip=6pt,
        afterskip=6pt
    },
    contentsname={\centerline{\normalfont\heiti\zihao{3}\texorpdfstring{目\qquad 录}{目录}}},
    bibname={\centerline{\heiti\zihao{3}参考文献}},
    paragraph={format={\zihao{-4}}}
}

% 研究生目录缩进：section 55pt = 1.5em(章节号) + 2.3em(标签宽) + 间距, subsection 70pt = 55pt + 15pt(层级偏移)
\dottedcontents{section}[55pt]{}{2.3em}{5pt}
\dottedcontents{subsection}[70pt]{}{2.3em}{5pt}
\fi

% 原始模板规范：本科tocdepth=2（目录显示到subsection），研究生tocdepth=4（目录显示到paragraph）
% 原始模板中研究生tocdepth在commands.tex中设为4，覆盖contents.tex中的2
\ifcauc@bachelor
\setcounter{tocdepth}{2}
\else
\setcounter{tocdepth}{4}
\fi

\renewcommand{\thesection}{\arabic{chapter}.\arabic{section}}
\renewcommand{\thesubsection}{\arabic{chapter}.\arabic{section}.\arabic{subsection}}
\renewcommand{\thesubsubsection}{\arabic{chapter}.\arabic{section}.\arabic{subsection}.\arabic{subsubsection}}

% 图/表/公式编号格式在 commands.tex 中 \counterwithin 之后设置，避免被覆盖

\ifcauc@bachelor
\IfFileExists{gbt7714-2005-numerical.bst}{%
  \bibliographystyle{gbt7714-2005-numerical}%
}{%
  \ClassWarning{cauc_thesis}{%
    Bibliography style file 'gbt7714-2005-numerical.bst' not found.%
    \MessageBreak
    Falling back to 'plain' style.%
    \MessageBreak
    Install: tlmgr install gbt7714%
  }%
  \bibliographystyle{plain}%
}%
\else
\IfFileExists{gbt7714-numerical.bst}{%
  \bibliographystyle{gbt7714-numerical}%
}{%
  \ClassWarning{cauc_thesis}{%
    Bibliography style file 'gbt7714-numerical.bst' not found.%
    \MessageBreak
    Falling back to 'plain' style.%
    \MessageBreak
    Install: tlmgr install gbt7714%
  }%
  \bibliographystyle{plain}%
}%
\fi
